\documentclass[11pt]{article}
\usepackage[a4paper,margin=1in]{geometry}
\usepackage{fontspec}
\usepackage[boldfont=false]{xeCJK}
\setCJKmainfont{FandolSong-Regular.otf}
\usepackage{amsmath}
\usepackage{amssymb}
\usepackage{array}
\usepackage{booktabs}
\usepackage{graphicx}
\usepackage{adjustbox}
\usepackage{enumitem}
\setlist[itemize]{topsep=2pt,partopsep=0pt,parsep=0pt,itemsep=2pt,leftmargin=1.4em}
\setlist[enumerate]{topsep=2pt,partopsep=0pt,parsep=0pt,itemsep=2pt,leftmargin=1.6em}
\usepackage{parskip}
\usepackage{fvextra}
\fvset{breaklines=true,breakanywhere=true,fontsize=\small}
\setlength{\parindent}{0pt}

\begin{document}

\begin{center}
  {\LARGE\bfseries An introduction to AS Level organic chemistry}\\[0.35em]
  {\large A Level Chemistry}
\end{center}
\vspace{0.6em}

\section*{Formulas, functional groups and naming}

A \textbf{hydrocarbon} 碳氢化合物 is a compound of only carbon and hydrogen. \textbf{Alkanes} 烷烃 are the simplest hydrocarbons and have no functional group.

A \textbf{functional group} 官能团 is the reactive part of a molecule. It decides the physical and chemical properties of the compound, so molecules with the same functional group behave alike (for example the $\text{–OH}$ group in alcohols).

\subsection*{Types of formula}

\begin{center}
\adjustbox{max width=\linewidth}{%
\begin{tabular}{ll}
\toprule
\textbf{Formula} & \textbf{What it shows} \\
\midrule
\textbf{general formula} 通式 & the pattern for a whole family, e.g. alkanes are $\text{C}_n\text{H}_{2n+2}$ \\
\textbf{molecular formula} 分子式 & the actual number of each atom, e.g. $\text{C}_4\text{H}_{10}$ \\
\textbf{structural formula} 结构式 & the groups in order, e.g. $\text{CH}_3\text{CH}_2\text{CH}_2\text{CH}_3$ \\
\textbf{displayed formula} 展开式 & every atom and every bond drawn out \\
\textbf{skeletal formula} 骨架式 & lines for bonds; carbons at corners, hydrogens on carbon not shown \\
\bottomrule
\end{tabular}}
\end{center}

You can read off the \textbf{empirical formula} 实验式 (simplest ratio) from any of these.

\subsection*{Naming}

Use systematic \textbf{nomenclature} 命名法: a stem for the number of carbons (meth-, eth-, prop-, but-, pent-, hex- for 1 to 6), an ending for the functional group, and numbers to show where groups are.

\section*{Characteristic organic reactions}

\subsection*{Some key terms}

\begin{itemize}
\item a \textbf{homologous series} 同系列 is a family of compounds with the same functional group and general formula, where each member differs by $\text{CH}_2$.

\item a \textbf{saturated} 饱和 compound has only single C–C bonds; an \textbf{unsaturated} 不饱和 compound has a C=C double bond (or a triple bond).

\end{itemize}

\subsection*{Breaking bonds}

A covalent bond can break in two ways:

\begin{itemize}
\item \textbf{homolytic fission} 均裂: the bond splits evenly, one electron to each atom. This makes two \textbf{free radicals} 自由基 (species with an unpaired electron).

\item \textbf{heterolytic fission} 异裂: the bond splits unevenly, both electrons going to one atom. This makes two ions.

\end{itemize}

\subsection*{Attacking species}

\begin{itemize}
\item a \textbf{nucleophile} 亲核试剂 is a species with a lone pair that is attracted to a positive (electron-poor) centre.

\item an \textbf{electrophile} 亲电试剂 is a species attracted to a negative (electron-rich) centre, such as a C=C double bond.

\end{itemize}

\subsection*{Types of reaction}

\begin{center}
\adjustbox{max width=\linewidth}{%
\begin{tabular}{ll}
\toprule
\textbf{Reaction} & \textbf{What happens} \\
\midrule
\textbf{addition} 加成 & two molecules join to make one \\
\textbf{substitution} 取代 & one atom or group is swapped for another \\
\textbf{elimination} 消去 & a small molecule is removed, making a double bond \\
\textbf{hydrolysis} 水解 & a molecule is split apart by water \\
\textbf{condensation} 缩合 & two molecules join and a small molecule (such as water) is lost \\
\bottomrule
\end{tabular}}
\end{center}

For organic redox, the symbol $[\text{O}]$ stands for one oxygen atom from an oxidising agent, and $[\text{H}]$ for one hydrogen atom from a reducing agent.

\subsection*{Types of mechanism}

A free-radical reaction happens in three steps: \textbf{initiation} 引发 (radicals are made), \textbf{propagation} 增长 (radicals react and make new radicals), and \textbf{termination} 终止 (two radicals join and stop the chain).

The main mechanisms you meet are \textbf{free-radical substitution} 自由基取代 (alkanes), \textbf{electrophilic addition} 亲电加成 (alkenes), \textbf{nucleophilic substitution} 亲核取代 (halogenoalkanes) and \textbf{nucleophilic addition} 亲核加成 (carbonyls). In mechanisms, a \textbf{curly arrow} 弯箭头 shows a pair of electrons moving; it starts at a bond or a \textbf{lone pair} 孤对电子.

\section*{Shapes of organic molecules}

Organic molecules can be straight-chained, branched or \textbf{cyclic} 环状 (in a ring).

The shape around a carbon depends on its \textbf{hybridisation} 杂化:

\begin{center}
\adjustbox{max width=\linewidth}{%
\begin{tabular}{llll}
\toprule
\textbf{Hybridisation} & \textbf{Bonds} & \textbf{Shape} & \textbf{Angle} \\
\midrule
$\text{sp}^3$ & 4 single & tetrahedral & $109.5°$ \\
$\text{sp}^2$ & 1 double + 2 single & \textbf{planar} 平面 (flat) & $120°$ \\
$\text{sp}$ & 1 triple (or 2 doubles) & linear & $180°$ \\
\bottomrule
\end{tabular}}
\end{center}

Every single bond is a \textbf{sigma bond} σ键, made by direct overlap. A double bond is one sigma bond plus one \textbf{pi bond} π键, made by sideways overlap of p orbitals. Ethene is planar because of its $\text{sp}^2$ carbons.

\section*{Isomerism}

\textbf{Isomers} are compounds with the same molecular formula but a different arrangement of atoms. This is called \textbf{isomerism} 异构.

\subsection*{Structural isomerism}

In \textbf{structural isomerism} 结构异构 the atoms are joined in a different order. There are three kinds:

\begin{itemize}
\item \textbf{chain isomerism} 链异构: the carbon chain is branched in different ways.

\item \textbf{positional isomerism} 位置异构: the functional group is on a different carbon.

\item \textbf{functional group isomerism} 官能团异构: the atoms form a different functional group (for example an alcohol and an ether).

\end{itemize}

\subsection*{Stereoisomerism}

In \textbf{stereoisomerism} 立体异构 the atoms are joined in the same order but point in different directions in space.

\begin{itemize}
\item \textbf{geometrical isomerism} 几何异构 (cis/trans) happens at a C=C double bond. The pi bond stops the two carbons rotating (\textbf{restricted rotation} 受限旋转), so groups are fixed on the same side (\textbf{cis} 顺式) or opposite sides (\textbf{trans} 反式).

\item \textbf{optical isomerism} 旋光异构 happens at a \textbf{chiral} 手性 carbon — a \textbf{chiral centre} 手性中心 is a carbon with four different groups attached. Such a carbon gives two mirror-image forms called \textbf{enantiomers} 对映体. A molecule may have more than one chiral centre.

\end{itemize}

From a molecular formula you can deduce the possible isomers by trying different chains, positions and functional groups.


\end{document}
